%=======================================================================
%  Cloud Computing -- 100 Day Research Assignment
%  Group G18 | Section A
%  Cutoff 1: Research Proposal
%
%  Compile with: pdfLaTeX  (Overleaf: Menu > Compiler > pdfLaTeX)
%  Bibliography: BibTeX    (references.bib must be in the same folder)
%=======================================================================

\documentclass[conference]{IEEEtran}
\IEEEoverridecommandlockouts

\usepackage{cite}
\usepackage{amsmath,amssymb,amsfonts}
\usepackage{graphicx}
\usepackage{booktabs}
\usepackage{array}
\usepackage{url}
\usepackage[hidelinks]{hyperref}
\usepackage{listings}
\usepackage{xcolor}

% --- listing style for YAML / shell snippets -------------------------
\lstdefinestyle{cfg}{
  basicstyle=\ttfamily\footnotesize,
  keywordstyle=\color{blue!60!black},
  commentstyle=\color{gray},
  breaklines=true,
  frame=single,
  framesep=4pt,
  showstringspaces=false
}
\lstset{style=cfg}

% --- convenience macros ----------------------------------------------
\newcommand{\rqbox}[1]{%
  \begin{center}
  \fbox{\begin{minipage}{0.94\columnwidth}\itshape #1\end{minipage}}
  \end{center}}

\begin{document}

\title{When Self-Healing Hurts: An Empirical Study of\\
Liveness-Probe-Induced Restart Cascades in\\
Kubernetes Under Overload}

\author{
\IEEEauthorblockN{M. Yasir Suhail}
\IEEEauthorblockA{\textit{Department of Computer Science}\\
\textit{FAST-NUCES, Islamabad}\\
Section A, Group G18}
\and
\IEEEauthorblockN{M. Aashir Sheraz}
\IEEEauthorblockA{\textit{Department of Computer Science}\\
\textit{FAST-NUCES, Islamabad}\\
Section A, Group G18}
}

\maketitle

%======================================================================
\begin{abstract}
Kubernetes restarts containers that fail their liveness probes, on the
assumption that a restart repairs a broken container and therefore
improves service availability. This assumption breaks down under
overload. When an application is merely slow rather than dead, its
health endpoint may exceed the probe timeout, the container is
restarted, and its share of the offered load is redistributed onto the
surviving replicas, pushing them toward the same threshold. The
resulting positive feedback loop is widely described in operational
practice, but it has not been characterised quantitatively. Existing
academic work on Kubernetes probes measures how quickly probes
\emph{detect} genuine failures; none measures the conditions under
which probes \emph{manufacture} failures. This proposal sets out a
controlled experimental study that sweeps the liveness-probe
configuration space under three load regimes on a multi-node cluster
and locates the boundary beyond which restart-based self-healing
reduces availability relative to running with no liveness probe at
all. The intended contributions are an empirically derived instability
boundary, a quantification of the availability penalty inside that
region, and a reproducible open-source testbed.
\end{abstract}

\begin{IEEEkeywords}
Kubernetes, self-adaptive systems, liveness probes, container
orchestration, service availability, cascading failure, fault
tolerance
\end{IEEEkeywords}

%======================================================================
\section{Group and Roles}

This research is conducted by Group~G18, Section~A. Both members
contribute to experimental design, analysis, and writing; the division
below reflects primary ownership rather than exclusive responsibility.

\begin{table}[h]
\centering
\caption{Group Composition and Role Allocation}
\label{tab:roles}
\footnotesize
\begin{tabular}{@{}p{0.26\columnwidth}p{0.16\columnwidth}p{0.44\columnwidth}@{}}
\toprule
\textbf{Member} & \textbf{Role} & \textbf{Primary Responsibility} \\
\midrule
M. Yasir Suhail & Primary /
Infrastructure & Cluster provisioning, service under test,
observability stack, experiment automation \\
\addlinespace
M. Aashir Sheraz & Co-investigator /
Analysis & Load modelling, statistical analysis, literature
review, manuscript preparation \\
\bottomrule
\end{tabular}
\end{table}

%======================================================================
\section{Research Area and Title}

\textbf{Research Area:} Autonomous / Self-Adaptive Cloud Computing.

\textbf{Title:} \emph{When Self-Healing Hurts: An Empirical Study of
Liveness-Probe-Induced Restart Cascades in Kubernetes Under Overload.}

The work sits squarely within self-adaptive cloud computing because
its object of study is an autonomic control loop --- the kubelet's
monitor--decide--act cycle over container health --- and the
conditions under which that loop becomes destabilising rather than
corrective.

%======================================================================
\section{Problem Statement and Motivation}

Container orchestrators are expected to detect and repair failures
without human intervention, and this self-healing capability is a
principal reason for their adoption \cite{casalicchio2019}. Kubernetes
implements it through two mechanisms: redundant replica Pods, and the
restarting of containers judged unhealthy. Health is judged by
\emph{probes}: the kubelet periodically issues an HTTP request, gRPC
call, or command against each container, and a container that fails a
configured number of consecutive liveness probes is restarted
\cite{roberts2025}.

The mechanism rests on an implicit assumption --- that a container
which fails to answer a probe in time is broken, and that restarting
it is therefore a repair. Under overload this assumption is unsound. A
saturated container is not broken; it is slow. If its health endpoint
is served by the same constrained resources as production traffic,
queueing delay inflates the probe response time past the configured
timeout. The kubelet then terminates a container that was still
serving useful work.

The second-order effect is the more serious one. The terminated
replica's share of the offered load is redistributed across the
remaining replicas, raising their utilisation and their health-endpoint
latency, which makes their own probes more likely to time out. The
system can enter a positive feedback loop in which the repair
mechanism is itself the source of continued failure. The Kubernetes
documentation warns that incorrect liveness-probe implementation can
cause cascading failures, and production incident reports describe the
phenomenon concretely --- one widely cited case required relaxing the
probe timeout and removing a CPU limit before the loop would break.

The motivation for studying this is practical. Probe parameters are
four integers in a YAML manifest, typically copied between projects
and rarely tuned against measured application behaviour. Practitioners
are advised to configure them ``carefully,'' but no published study
tells them where the boundary of safety lies, how far inside the
dangerous region a given configuration sits, or how much availability
is lost when it does. A misconfiguration that is harmless at 50\%
utilisation may be catastrophic at 150\%, and nothing in the current
literature allows an operator to anticipate that transition.

%======================================================================
\section{Research Gap}

Prior work on Kubernetes health monitoring addresses adjacent
questions but leaves this one open.

\textbf{Probe speed, not probe correctness.} Roberts et
al.~\cite{roberts2025} model and measure the time Kubernetes probes
take to detect genuine container failure, and propose signal-based
monitoring as a faster alternative. Probes are their baseline, not
their subject. Their evaluation deploys a single Pod per service on a
single node with throughput capped at a fixed rate, and injects
failure by terminating a process once at a fixed time. With one
replica per service there is no load redistribution, so the feedback
loop studied here cannot form by construction. The authors themselves
note that their failure model is simple and does not include bursts or
randomised arrival.

\textbf{An unexplained observation in that same work.} Under their
aggressively tuned ``fast probe'' configuration, the container
restarted several times on average and service availability fell by
approximately 4\% relative to the signal-based alternative, with the
additional restarts attributed to failing liveness probes. The authors
state that it remains unclear why most of these restarts occur after
the injected failure. The phenomenon was observed as a side effect and
left unexplained.

\textbf{An explicitly stated open problem.} The same paper's future
work proposes injecting network latency to emulate slow HTTP responses
and states the expectation that poll-based monitoring will erroneously
detect high latency as container failure. The prediction is on record;
the experiment has not been performed.

\textbf{The opposite failure mode has been studied.} Flora et
al.~\cite{flora2022} find that Kubernetes probes fail to detect
certain classes of fault, including business-logic errors and long
response times --- that is, false negatives. The complementary case,
in which probes fire on healthy-but-loaded containers, has not
received equivalent treatment.

Broader microservice fault taxonomies \cite{silva2022,zhou2021} and
resilience-testing frameworks \cite{heorhiadi2016} catalogue
application- and infrastructure-level faults, but do not treat the
orchestrator's own health-check subsystem as a fault source.

\textbf{Statement of the gap.} To the best of our knowledge, no
published study systematically varies Kubernetes liveness-probe
parameters under sustained overload, with replica counts sufficient
for load redistribution, and measures the resulting effect on service
availability against a no-probe control. The failure mode is
established as operational folklore and observed incidentally in the
literature, but it has not been characterised.

%======================================================================
\section{Research Question}

\rqbox{Under sustained overload, at what combinations of Kubernetes
liveness-probe parameters (\texttt{timeoutSeconds},
\texttt{periodSeconds}, \texttt{failureThreshold}) and container CPU
limits does restart-based self-healing reduce service availability
rather than restore it, and how large is the resulting degradation
relative to a deployment configured with no liveness probe?}

Two subsidiary questions follow:

\begin{itemize}
\item \textbf{SQ1.} How does offered load mediate the relationship
between probe aggressiveness and restart frequency --- that is, does
the instability boundary shift with utilisation, and by how much?
\item \textbf{SQ2.} Does the severity of degradation scale with
replica count, as the load-redistribution mechanism would predict?
\end{itemize}

%======================================================================
\section{Objectives}

\begin{enumerate}
\item \textbf{O1 --- Construct a reproducible testbed.} Build a
multi-node Kubernetes cluster hosting an instrumented service whose
health endpoint contends with production traffic for the same
constrained resources, together with a deterministic load generator
and a metrics pipeline capable of attributing each restart to a probe
failure.

\item \textbf{O2 --- Characterise the load--latency--probe coupling.}
Establish empirically how offered load translates into health-endpoint
response latency, and identify the utilisation level at which that
latency begins to approach representative probe timeouts. This
establishes the causal chain on which the remainder of the study
depends.

\item \textbf{O3 --- Locate the instability boundary.} Sweep the probe
parameter space across three load regimes and identify the region in
which restart cascades emerge, using a no-liveness-probe deployment as
the control arm.

\item \textbf{O4 --- Quantify the penalty and derive guidance.}
Measure the availability, error-rate, and tail-latency cost of
operating inside the unstable region, and translate the findings into
a defensible configuration envelope for practitioners.
\end{enumerate}

%======================================================================
\section{Proposed Methodology}

The study is a controlled factorial experiment. The strategy is to fix
every source of variation except probe configuration and offered load,
so that any difference in availability is attributable to the health-check
subsystem rather than to the application or the autoscaler.

\subsection{Experimental Platform}

A \texttt{kind} (Kubernetes-in-Docker) cluster of one control-plane
node and three worker nodes, running on a single commodity laptop.
Container CPU limits are set low (on the order of 200~millicores) so
that the cgroup quota binds well before the host's own capacity does;
this keeps the measured saturation point a property of the pod
configuration rather than of the test machine. Horizontal Pod
Autoscaling is disabled throughout and replica count is held fixed, so
that the only adaptive mechanism active in the system is the one under
study.

\subsection{Service Under Test}

A purpose-built HTTP service exposing two endpoints:

\begin{itemize}
\item \texttt{/work} --- performs bounded CPU-intensive work,
representing production traffic.
\item \texttt{/healthz} --- the probe target, served from the same
bounded worker pool as \texttt{/work}.
\end{itemize}

The shared worker pool is the essential design decision. It reproduces
the common real-world pattern in which a health endpoint is handled by
the same thread pool, event loop, or CPU quota as application
requests, and it is the mechanism by which load is transmitted into
probe latency. A service whose health endpoint is independent of load
cannot exhibit the phenomenon under study.

\subsection{Fault and Workload Model}

No faults are injected. This is a deliberate departure from prior work
and is central to the contribution: the study asks whether the
orchestrator generates failures in a system where none exist. The
independent stressor is offered load alone, applied as a
constant-arrival-rate workload at three levels --- approximately
50\%, 100\%, and 150\% of measured single-replica saturation
throughput --- using \texttt{k6}.

\subsection{Experimental Design}

\begin{table}[h]
\centering
\caption{Factors and Levels}
\label{tab:factors}
\footnotesize
\begin{tabular}{@{}p{0.34\columnwidth}p{0.30\columnwidth}p{0.22\columnwidth}@{}}
\toprule
\textbf{Factor} & \textbf{Levels} & \textbf{Type} \\
\midrule
\texttt{timeoutSeconds} & 1, 2, 5 & Independent \\
\texttt{failureThreshold} & 1, 3, 10 & Independent \\
Offered load & 50\%, 100\%, 150\% & Independent \\
\texttt{periodSeconds} & 5 (fixed) & Controlled \\
Replica count & 6 (fixed) & Controlled \\
CPU limit & 200m (fixed) & Controlled \\
Liveness probe absent & control arm & Baseline \\
\bottomrule
\end{tabular}
\end{table}

The full factorial gives $3 \times 3 \times 3 = 27$ treatment
configurations, plus three control runs (one per load level) in which
no liveness probe is defined. Each configuration is executed five
times to accommodate the stochastic nature of cascade onset, giving
150 runs of approximately ten minutes each. Execution is fully
scripted and unattended.

\subsection{Instrumentation and Analysis}

Prometheus with \texttt{kube-state-metrics} records restart counts and
Pod readiness over time; kubelet events attribute each restart to a
specific probe failure; \texttt{k6} reports client-observed success
rate and latency distribution. Treatment arms are compared against the
matched control using non-parametric tests, and the parameter space is
visualised as a heatmap over which the instability boundary is drawn.

\subsection{Threats to Validity}

Single-host deployment constrains achievable scale, and a synthetic
single-service workload is not a substitute for a production
microservice graph. Both are acknowledged limitations; the design
compensates by holding the confounded variables fixed and by reporting
the saturation point as a measured rather than assumed quantity.
Should the cascade prove difficult to induce reproducibly, the
mitigation is to tighten the CPU limit and reduce the probe timeout
until the coupling established in O2 is strong enough to trip the
mechanism.

%======================================================================
\section{Expected Metrics and Outcomes}

\subsection{Metrics}

\begin{table}[h]
\centering
\caption{Measured Quantities}
\label{tab:metrics}
\footnotesize
\begin{tabular}{@{}p{0.30\columnwidth}p{0.60\columnwidth}@{}}
\toprule
\textbf{Metric} & \textbf{Definition} \\
\midrule
Restart count & Cumulative container restarts attributable to
liveness-probe failure \\
\addlinespace
Availability & Time-weighted fraction of the run during which the
service sustained target throughput \\
\addlinespace
Success rate & Proportion of client requests completing without
timeout or HTTP 5xx \\
\addlinespace
Tail latency & P95 and P99 client-observed response time \\
\addlinespace
Recovery time & Interval from first probe-induced restart to
restoration of steady-state throughput \\
\addlinespace
Cascade depth & Maximum number of replicas simultaneously
unavailable \\
\bottomrule
\end{tabular}
\end{table}

\subsection{Anticipated Outcomes}

\begin{enumerate}
\item \textbf{An instability boundary.} A demarcation in
$(\texttt{timeout}, \texttt{failureThreshold}, \texttt{load})$ space
separating configurations in which the liveness probe improves
availability from those in which it degrades it.

\item \textbf{A quantified penalty.} The magnitude of availability
loss inside the unstable region, expressed relative to the no-probe
control. The central and most counter-intuitive expected finding is
the existence of a non-empty region in which \emph{no liveness probe
is strictly preferable to a misconfigured one}.

\item \textbf{A configuration envelope.} Evidence-based guidance
relating safe probe parameters to measured application saturation
behaviour, replacing the qualitative advice presently available.

\item \textbf{A reproducible artifact.} A public repository
containing the service, manifests, automation, raw results, and
analysis notebooks.
\end{enumerate}

A negative result --- that cascades do not materialise within the
swept range --- remains a publishable finding, since it would bound
the practical severity of a failure mode currently treated as a
serious operational risk.

%======================================================================
\section{Project Resources}

\begin{table}[h]
\centering
\caption{Deliverable Infrastructure}
\label{tab:resources}
\footnotesize
\begin{tabular}{@{}p{0.24\columnwidth}p{0.66\columnwidth}@{}}
\toprule
\textbf{Resource} & \textbf{Status} \\
\midrule
Overleaf & Project created; read link shared with instructor \\
GitHub & Public repository initialised \\
Google Sheet & Row for G18 completed \\
\bottomrule
\end{tabular}
\end{table}

%======================================================================
\bibliographystyle{IEEEtran}
\bibliography{references}

\end{document}
